% !TEX program = lualatex
%!TEX TS-program = lualatex
\documentclass[12pt,a4paper]{article}
\usepackage{fontspec}
\usepackage{polyglossia}
\setdefaultlanguage{polish}
\usepackage{geometry}
\geometry{margin=2.1cm}
\usepackage[table]{xcolor}
\usepackage{longtable}
\usepackage{array}
\usepackage{booktabs}
\usepackage{enumitem}
\usepackage{ragged2e}
\definecolor{tableHeader}{HTML}{E8EEF7}
\definecolor{tableStripe}{HTML}{F7F9FC}
\newcolumntype{L}[1]{>{\RaggedRight\arraybackslash}p{#1}}
\newcolumntype{C}[1]{>{\Centering\arraybackslash}p{#1}}
\renewcommand{\arraystretch}{1.22}
\setlength{\tabcolsep}{4pt}
\setlist[itemize]{leftmargin=*, itemsep=2pt, topsep=2pt}
\setlength{\parindent}{0pt}
\setlength{\parskip}{6pt}
\emergencystretch=3em
\sloppy
\begin{document}

\section*{Raport ze sprintu nr 1 - Calisia Banking App}
Data raportu: 03.05.2026

Okres sprintu: 25.04 - 30.04.2026

Product Owner / Scrum Master: Jakub Wołczko

\subsection*{Lista osób}
\begin{itemize}
\item Kubryn Jeremiasz
\item Łaszuk Jakub
\item Makarevich Ksenia
\item Nuckowski Michał
\item Ruczyńska Nikola
\item Wołczko Jakub
\item Olszewski Dawid
\end{itemize}

\subsection*{Product Owner i Scrum Master sprintu}
Jakub Wołczko

\subsection*{Lista historyjek na sprint}
\begingroup
\footnotesize
\rowcolors{2}{tableStripe}{white}
\begin{longtable}{@{}C{0.08\textwidth}L{0.62\textwidth}C{0.10\textwidth}L{0.12\textwidth}@{}}
\toprule
\rowcolor{tableHeader}
\textbf{Kod} & \textbf{Historyjka} & \textbf{SP} & \textbf{Status} \\
\midrule
\endfirsthead
\toprule
\rowcolor{tableHeader}
\textbf{Kod} & \textbf{Historyjka} & \textbf{SP} & \textbf{Status} \\
\midrule
\endhead
S1-01 & Przygotowanie solucji backendowej z warstwami Api, Application, Domain, Infrastructure, Contracts i ReadModel & 8 & Done \\
S1-02 & Wdrożenie wzorca CQRS dla komend i zapytań & 8 & Done \\
S1-03 & Konfiguracja projektu frontendowego React/Vite z routingiem, stylami globalnymi i layoutami & 5 & Done \\
S1-04 & Przygotowanie skryptów baz CalisiaWriteDb i CalisiaReadDb & 5 & Done \\
S1-05 & Dodanie dokumentacji uruchomieniowej w README & 3 & Done \\
S1-06 & Przygotowanie kolekcji Bruno dla podstawowych endpointów & 3 & Done \\
\bottomrule
\end{longtable}
\endgroup

\subsection*{Podsumowanie sprintu}
Cel sprintu: Fundamenty architektury i środowiska. Przygotowanie technicznego fundamentu aplikacji, aby zespół mógł rozwijać backend, frontend i bazy danych w jednym spójnym projekcie.

Zakres sprintu obejmował 6 historyjek o łącznej estymacji 32 SP. Wszystkie historyjki zostały dowiezione ze statusem Done, dlatego osiągnięta prędkość sprintu wyniosła 32 SP.

Przygotowano kompletną bazę techniczną projektu: wielowarstwową solucję backendową, podstawowe abstrakcje CQRS, projekt frontendowy React/Vite, skrypty baz danych, dokumentację uruchomieniową oraz początkową kolekcję Bruno do weryfikacji API.

Sprint stworzył stabilny fundament pod dalszą implementację funkcji biznesowych. Dzięki temu kolejne prace mogą być prowadzone równolegle w backendzie, frontendzie, bazach danych i testach bez konieczności ciągłego zmieniania struktury projektu.

\subsubsection*{Najważniejsze rezultaty}
\begin{itemize}
\item Utworzono bazową strukturę aplikacji i repozytorium.
\item Wprowadzono wzorzec CQRS jako standard implementacji komend i zapytań.
\item Przygotowano frontend React/Vite z routingiem oraz podstawowymi layoutami.
\item Dodano skrypty baz danych i pierwszą kolekcję Bruno.
\end{itemize}

\subsection*{Retrospekcja sprintu}
Retrospekcja pokazała, że zespół dobrze rozpoczął projekt, ale potrzebuje mocniej spisywać konwencje techniczne już na początku prac.

\subsubsection*{Co się udało}
\begin{itemize}
\item Dobre rozdzielenie odpowiedzialności pomiędzy warstwami backendu.
\item Szybkie uzgodnienie wspólnej struktury repozytorium.
\item Konsekwentne zastosowanie CQRS od początku projektu.
\item Wczesne dodanie kolekcji Bruno jako narzędzia do testowania API.
\end{itemize}
\subsubsection*{Poziom energii zespołu}
Średni poziom energii zespołu: 4,1/5.

Energia zespołu była pozytywna i produktywna. Jasny cel sprintu oraz widoczny przyrost techniczny zwiększyły poczucie kontroli nad projektem.

\subsubsection*{Czego udało się nauczyć}
\begin{itemize}
\item Zespół utrwalił zasady pracy w architekturze warstwowej.
\item Warto spisywać konwencje projektowe przed rozwojem kolejnych funkcji.
\item Dokumentacja uruchomieniowa i Bruno ułatwiają onboarding oraz testowanie.
\end{itemize}
\subsubsection*{Wnioski i działania na kolejny sprint}
\begin{itemize}
\item Spisać listę konwencji projektowych dla nazw folderów, namespaceów i endpointów.
\item Dodać do README krótką sekcję szybkiej weryfikacji środowiska.
\item Aktualizować kolekcję Bruno razem z rozwojem endpointów.
\end{itemize}

\subsection*{Następny sprint}
Product Owner / Scrum Master w następnym sprincie: Jakub Łaszuk

Główny cel następnego sprintu: Klient, rejestracja i bezpieczne hasła. Umożliwienie klientowi założenia profilu oraz przygotowanie bezpiecznego przechowywania danych logowania.

\end{document}
